\documentclass{tufte-handout}

\usepackage{style}

\begin{document}

\tma{04} % TMA number

\section{Part A}

%Question 4
\begin{question}

    Let
    \[ \bm{F}  = \Z_{2}(\omega) \]

    Where,
    \[ \omega^{4} = \omega + 1 \]

    \qpart

    For \( z \in F \) we have \( z + z = 0 \) because \( z + z = 2z \equiv 0 \pmod{2} \), so \( z = -z \) for all \( z \in F \).



\qpart

\qsubpart

    To find the multiplicative inverse of \( 1 + \omega^{2} \in F \) we by an arbitary element of \( K \).

    \begin{align*}
    \stext{The given elemnt can be wriiten as,}
    a\omega^{3} + b\omega^{2} + c\omega + d
    \stext{Where \( a,b,c,d \in \Z_{2} \), so we need to find \( a,b,c,d \) such that,}
    1 &= (\omega^{2} + 1)(a\omega^{3} + b\omega^{2} + c\omega + d)\\[8pt]
    &= a\omega^{5} + b\omega^{4} + c\omega^{3} + d\omega^{2} + a\omega^{3} + b\omega^{2} + c\omega + d\\[8pt]
    \stext{Since \( \omega^{4} = \omega + 1 \),}
    \stext{ we can write \( \omega^{5} = \omega\omega^{4} = \omega(\omega + 1) = \omega^{2} + \omega \), so,}
    &= a(\omega^{2} + \omega) + b(\omega + 1) + c\omega^{3} + d\omega^{2} + a\omega^{3} + b\omega^{2} + c\omega + d\\[8pt]
    &= (b + d) + (a + b + c)\omega + (a + b + d)\omega^{2} + (a + c)\omega^{3}\\[8pt]
    \end{align*}

\vspace{1cm}

\qsubpart

This gives us the system of equations,
\begin{align*}
b + d &= 1 \in \Z_{2}\\[8pt]
a + b + c &= 0 \in \Z_{2}\\[8pt]
a + b + d &= 0 \in \Z_{2}\\[8pt]
a + c &= 0 \in \Z_{2}\\[8pt]
\end{align*}

From the last equation we have \( c = a \), and substituting this into the second equation gives us \( b = 0 \), and substituting this into the third equation 
gives us \( d = a \), and substituting this into the first equation gives us \( a = 1 \).

Hence, the multiplicative inverse of \( 1 + \omega^{2} \) is,

\[ \omega^{3} + \omega + 1 \]

\vspace{2cm}

\qpart


\end{question}

\end{document}